\begin{statementbox}
\textbf{C087｜双曲线焦点弦长公式与取值范围}

设横轴型双曲线
\[
H:\ \frac{x^2}{a^2}-\frac{y^2}{b^2}=1,\quad
 a>0,\ b>0,\quad c=\sqrt{a^2+b^2},\quad F_2(c,0).
\]
直线 $\ell$ 过焦点 $F_2$，与双曲线交于两个\textbf{不同的有限点} $A,B$，
线段 $AB$ 为\textbf{焦点弦}。两焦点中的另一个为 $F_1(-c,0)$；
对 $F_1$ 完全对称，结论不变。
\begin{center}% 精确坐标：H: x^2/9-y^2/16=1, c=5, b^2/a=16/3.
% 两图中蓝色曲线由解析式 x=+-3 sqrt(1+y^2/16) 绘制，非手绘。
\begin{tikzpicture}[x=0.39cm,y=0.39cm,>=stealth, font=\small]
\begin{scope}
  \draw[->,gray] (-6.8,0)--(7.3,0) node[right]{$x$};
  \draw[->,gray] (0,-7.8)--(0,7.8) node[above]{$y$};
  \draw[thick,blue!75!black,domain=-7:7,samples=90] plot ({3*sqrt(1+\x*\x/16)},{\x});
  \draw[thick,blue!75!black,domain=-7:7,samples=90] plot ({-3*sqrt(1+\x*\x/16)},{\x});
  \draw[very thick,red!75!black] (5,-5.33333333)--(5,5.33333333);
  \fill (5,0) circle (2.1pt) node[below right]{$F_2$};
  \fill[red!75!black] (5,5.33333333) circle (1.8pt) node[above left=5pt]{$A$};
  \fill[red!75!black] (5,-5.33333333) circle (1.8pt) node[below left=5pt]{$B$};
  \node[align=center] at (0,-9.0) {同支焦点弦\quad $x=c$\\最短 $2b^2/a$};
\end{scope}
\begin{scope}[xshift=6.6cm]
  \draw[->,gray] (-6.8,0)--(7.3,0) node[right]{$x$};
  \draw[->,gray] (0,-7.8)--(0,7.8) node[above]{$y$};
  \draw[thick,blue!75!black,domain=-7:7,samples=90] plot ({3*sqrt(1+\x*\x/16)},{\x});
  \draw[thick,blue!75!black,domain=-7:7,samples=90] plot ({-3*sqrt(1+\x*\x/16)},{\x});
  \draw[very thick,red!75!black] (-3,0)--(3,0);
  \fill (5,0) circle (2.1pt) node[above right]{$F_2$};
  \fill[red!75!black] (-3,0) circle (1.8pt) node[above left=5pt]{$A$};
  \fill[red!75!black] (3,0) circle (1.8pt) node[above right=5pt]{$B$};
  \node[align=center] at (0,-9.0) {异支焦点弦\quad $y=0$\\最短 $2a$};
\end{scope}
\end{tikzpicture}
\end{center}

\textbf{结论一（斜率形式，仅针对非竖直直线）：}
令 $\ell:y=k(x-c)$，若 $k\ne\pm\dfrac ba$，则
\[
\boxed{|AB|=\frac{2ab^2(1+k^2)}{|a^2k^2-b^2|}}.
\]
\textbf{结论二（倾斜角形式，包含竖直线）：}
若 $\ell$ 的倾斜角为 $\theta\in[0,\pi)$，则
\[
\boxed{|AB|=\frac{2ab^2}{|a^2-c^2\cos^2\theta|}},
\quad \cos\theta\ne\pm\frac ac.
\]
特别地，$\theta=\dfrac\pi2$（直线 $x=c$）合法，弦长为 $2b^2/a$。

\end{statementbox}
\clearpage
\begin{statementbox}
\textbf{结论三（分类、范围与取等）：}
\[
\boxed{\begin{array}{lcl}
\text{异支焦点弦：}& |k|<b/a,& |AB|\in[2a,+\infty),\\[3pt]
\text{同支焦点弦：}& |k|>b/a\ \text{或竖直},& |AB|\in[2b^2/a,+\infty).
\end{array}}
\]
异支在 $\ell:y=0$ 取最小值 $2a$；
同支在 $\ell:x=c$ 取最小值 $2b^2/a$。
\begin{center}% Actual L(theta), a=3,b=4,c=5; horizontal degree/10, vertical .14L.
\begin{tikzpicture}[font=\small,>=stealth]
\draw[gray!45,rounded corners=2pt] (0,0) rectangle (14.2,9.2);
\node at (7.1,8.55) {转动直线看弦长：临界处无界增长，水平和竖直分别取最小值};
\node at (7.1,7.75) {图取 $a=3,b=4,c=5$，只看 $0\le\theta\le\pi/2$};
\begin{scope}[shift={(2.3,1.6)}]
\draw[gray,->] (0,0)--(9.7,0) node[right] {$\theta$};
\draw[gray,->] (0,0)--(0,5.8) node[above] {$L$};
\draw[gray,dashed] ({acos(3/5)/10},0)--({acos(3/5)/10},5.8);
\draw[red!75!black,very thick,->,domain=0:{acos(sqrt((9+96/40)/25))},samples=100,variable=\v]
 plot ({\v/10},{.14*96/(25*cos(\v)^2-9)});
\draw[purple!80!black,very thick,<-,domain={acos(sqrt((9-96/40)/25))}:90,samples=100,variable=\v]
 plot ({\v/10},{.14*96/(9-25*cos(\v)^2)});
\fill[red!75!black] (0,.84) circle(2.2pt) node[left=6pt] {$2a=6$};
\fill[purple!80!black] (9,{.14*32/3}) circle(2.2pt) node[above right=5pt] {$2b^2/a=32/3$};
\node[below=5pt] at (0,0) {$0$};
\node[below=5pt] at (9,0) {$\pi/2$};
\node[below=5pt] at ({acos(3/5)/10},0) {$\theta_0=\arccos(3/5)$};
\node[red!75!black] at (2,2.8) {异支};
\node[purple!80!black] at (7.4,3.4) {同支};
\end{scope}
\node at (7.1,.45) {临界角 $\theta_0$ 对应渐近线方向；实心点说明两个最小值都能取到。};
\end{tikzpicture}
\end{center}
\textbf{边界：}$k=\pm b/a$ 时 $\ell$ 平行于渐近线 $y=\pm(b/a)x$，
与双曲线\textbf{只有一个有限交点}，不构成通常意义下的焦点弦；
$k$ 逼近该边界时弦长 $\to+\infty$。
\end{statementbox}
